\documentclass[a4paper,twoside]{article}

\usepackage{epsfig}
\usepackage{subcaption}
\usepackage{calc}
\usepackage{amssymb}
\usepackage{amstext}
\usepackage{amsmath}
\usepackage{amsthm}
\usepackage{multicol}
\usepackage{pslatex}
\usepackage{apalike}
\usepackage{color}
\usepackage{algorithm2e}
\usepackage[bottom]{footmisc}
\usepackage{CJKutf8}        
\usepackage{SCITEPRESS}     
\usepackage{comment}
\newtheorem{theorem}{Theorem}
\newtheorem{lemma}{Lemma}
\newtheorem{proposition}{Proposition}
\newtheorem{corollary}{Corollary}
\theoremstyle{definition}
\newtheorem{definition}{Definition}
\newtheorem{example}{Example}

\newcommand{\zh}[1]{\textcolor{red}{\textbf{Zhaohong says:} #1}}
\newcommand{\zhnew}[1]{\textcolor{red}{#1}}
\newcommand{\topc}{\mathrm{top}}

\newenvironment{jproof}{\begin{proof}\begin{CJK}{UTF8}{min}}{\end{CJK}\end{proof}}
\newcommand{\aidraft}{\footnote{The text of this section was drafted with a generative AI tool \cite{Anthropic26a} and checked by the authors; see the Acknowledgements.}}
\newcommand{\aiedit}{\footnote{The wording of this section was edited with a generative AI tool \cite{Anthropic26a} based on the authors' technical notes; see the Acknowledgements.}}

\begin{document}

\title{Reserve Systems with Score Feasibility: Reserved Category Protection, De-Reservation, and Strategy-Proofness}


\author{\authorname{Anonymous Author(s)}
\affiliation{Affiliation omitted for double-blind review}
\email{anonymous@example.org}
}

\keywords{Matching, Market Design, Affirmative Action, Reserve System, Strategy-Proofness, Serial Dictatorship.}

\abstract{Reserve systems are a standard tool for affirmative action in admissions: some seats are set aside for reserved categories, and the remaining seats are open to all students. We study a single-college model in which students are ranked by exam scores, and matchings must be score-feasible: a student can be admitted through a reserved category only if her score is at most a category-specific margin below that of the highest-scoring student rejected from the open category. This requirement is motivated by affirmative action in India, where the Supreme Court directed the government to consider bounding the gap between the cutoff scores of the open category and a reserved category. When quotas are fixed, a score-based serial dictatorship mechanism (SSD) is feasible, score-feasible, fair, Pareto efficient, and strategy-proof. When unfilled reserved seats can be transferred to the open category, we introduce reserved category protection and show that strategy-proofness, reserved category protection, and weak non-wastefulness cannot be satisfied simultaneously together with feasibility, score-feasibility, and fairness. In contrast, each pair of these three properties is achieved by a variant of SSD. Thus, we obtain a tight boundary between possibility and impossibility.}

\onecolumn \maketitle \normalsize \setcounter{footnote}{0} \vfill

\section{\uppercase{Introduction}}
\label{sec:introduction}

Matching theory has been successfully applied to admission problems such as school
choice and college admissions
\cite{GaSh62a,RoSo90a,BaSo99a,AbSo03b}.
In many of these markets, the designer must implement affirmative action.
A common instrument is a \emph{reserve system}: some seats are set aside for students
who belong to protected groups, while the remaining seats are open to all students.
Reserve systems are used, for example, in school choice programs in the United States
\cite{DKPS18a,DPS20a}, in university admissions in Brazil \cite{AyBo21a}, and, on a
very large scale, in admissions to public universities and in the allocation of
government jobs in India \cite{BCC+19b,SoYe22a,AyTu23a}.

When students are ranked by an entrance exam score, a reserve system inevitably
creates a tension between meritocracy and diversity \cite{imamura2025meritocracy}.
Two practical concerns arise.
First, a student admitted through a reserved category can have a much lower score than
a student rejected from the open category.
Second, reserved seats can remain unfilled while students who want open seats are
rejected.
Both concerns appear explicitly in Indian affirmative action.
A 2008 judgment of the Supreme Court of India (Ashoka Kumar Thakur vs.\ Union of India)
upheld a 27\% reservation for Other Backward Classes (OBC).
Under this judgment, unfilled OBC seats are to revert to the general category, and the
government is directed to set up a committee to consider setting the OBC cutoff at no
more than 10 marks below that of the general category
\cite{AyTu23a}.
The judgment does not specify a procedure to implement these directives
\cite{AyTu23a}.

In this paper, we study these two requirements in a stylized model with a single
college. The college has an open category and several reserved categories.
Each student has an exam score, is eligible for the open category and possibly for some
reserved categories, and has a strict preference over her eligible categories and an
outside option.
That is, as in models of matching with contracts \cite{HaMi05a,AyTu20a} and controlled
school choice with overlapping types \cite{KHIY17a}, a student can care about the
category under which she is admitted.
We formalize the bounded-gap requirement as \emph{score-feasibility}: a student can be
admitted under a reserved category $t$ only if her score is at least $e_o-\delta_t$, where
$e_o$ is the highest score among rejected students who find the open category
acceptable, and $\delta_t\ge 0$ is a margin given for each reserved category
(i.e., different categories can have different margins). 
Thus, $e_o$ plays the role of the cutoff score of the open category.
Our contributions are as follows.
\begin{comment}
    

\begin{itemize}
\item \emph{Fixed quotas.} When each category has a fixed quota, we show that the
score-based serial dictatorship mechanism (SSD) is feasible, score-feasible, fair,
Pareto efficient, and strategy-proof (Theorem~\ref{thm:ssd}).
\item \emph{Flexible quotas and an impossibility.} When unfilled reserved
seats can be used for the open category, the designer must decide how many reserved
seats to protect. We introduce \emph{reserved category protection} (RCP), where the
protection level is the number of students whose top choice is the reserved category.
We show that, in some cases, no matching satisfies feasibility, score-feasibility,
non-wastefulness, and RCP (Theorem~\ref{thm:imp1}).
Moreover, even if non-wastefulness is weakened to \emph{weak non-wastefulness} (WNW),
no fair mechanism satisfies strategy-proofness (SP), RCP, and WNW simultaneously
(Theorem~\ref{thm:imp2}).
\item \emph{A tight boundary.} We show that each pair of SP, RCP, and WNW is achievable
together with feasibility, score-feasibility, and fairness.
Reserved category protecting SSD satisfies RCP and WNW (Theorem~\ref{thm:rcpssd}),
flexible SSD satisfies SP and non-wastefulness (Theorem~\ref{thm:flexible}), and
fixed quota SSD, which simply runs SSD with the open quota
$q_o=Q-\sum_{t\in T}q_t$, satisfies SP and RCP (Theorem~\ref{thm:fqssd}).
Together with Theorem~\ref{thm:imp2}, these results give a tight boundary between
possibility and impossibility (Corollary~\ref{cor:boundary}).
\end{itemize}
\end{comment}

\begin{itemize}
\item \emph{Score-feasibility and fixed quotas.}
We formulate score-feasibility using the highest score among
rejected students who find the open category acceptable, rather
than the lowest score among admitted open-category students.
We establish a key structural property: removing the lowest-score
admitted student preserves score-feasibility.
Building on this property, we show that when each category has
a fixed quota, the score-based serial dictatorship mechanism (SSD)
is feasible, score-feasible, fair, Pareto efficient, and
strategy-proof (Theorem~\ref{thm:ssd}).
\item \emph{Flexible quotas and an impossibility.}
When unfilled reserved seats can be transferred to the open
category, the designer must decide how much reserve capacity
to protect.
We introduce \emph{reserved category protection} (RCP),
which protects reserve capacity according to students'
reported first-choice demand.
We show that, in some instances, no matching simultaneously
satisfies feasibility, score-feasibility, non-wastefulness,
and RCP (Theorem~\ref{thm:imp1}).
Moreover, even if non-wastefulness is weakened to
\emph{weak non-wastefulness} (WNW), no mechanism
can simultaneously satisfy feasibility, score-feasibility,
fairness, strategy-proofness (SP), RCP, and WNW
(Theorem~\ref{thm:imp2}).
\item \emph{A tight boundary through different protection levels.}
We analyze three variants of SSD corresponding to different
levels of reserve protection: no protection, demand-based
protection, and full quota protection.
We show that each pair of SP, RCP, and WNW is achievable
together with feasibility, score-feasibility, and fairness.
Reserved category protecting SSD (RCP-SSD) satisfies RCP and WNW
(Theorem~\ref{thm:rcpssd}),
flexible SSD satisfies SP and non-wastefulness
(Theorem~\ref{thm:flexible}), and
fixed quota SSD, which runs SSD with
$q_o=Q-\sum_{t\in T}q_t$, satisfies SP and RCP
(Theorem~\ref{thm:fqssd}).
Together with Theorem~\ref{thm:imp2}, these results establish
a tight boundary between possibility and impossibility.
\end{itemize}

The rest of this paper is organized as follows.
Section~\ref{sec:related} discusses related work.
Section~\ref{sec:model} introduces the model.
Section~\ref{sec:fixed} analyzes the fixed quota model, and
Section~\ref{sec:flexible} analyzes the flexible quota model.
Section~\ref{sec:boundary} presents the three mechanisms and the tight boundary.
Section~\ref{sec:conclusion} concludes.

\section{\uppercase{Related Work}}
\label{sec:related}



\subsection{Affirmative Action in Matching}

College admissions and school choice with affirmative action have been studied
extensively \cite{Abdu05a,AbSo03b}.
It is shown that quota-based affirmative action can hurt students, even though it is intended to help them \cite{Koji12a}.
This result motivated reserve-based approaches \cite{HYY13a} and the distinction between hard
bounds and soft bounds \cite{EHYY14a}.
Choice rules that balance priorities and diversity have been characterized axiomatically
\cite{EcYe15a,imamura2025meritocracy}, and slot-specific priorities provide a general
framework for reserve systems \cite{KoSo16a}.
The order in which reserved seats and open seats are processed is known to
significantly affect the outcome \cite{DKPS18a,DPS20a,PRS23a}.

Reserve systems with overlapping categories have been studied in
\cite{SoYe22a,AzSu25a}, and reserve systems are also used in healthcare
rationing \cite{PSUY24a,AzBr24a}.
In these rationing applications, the designer decides the category through which each
agent is served, whereas in our model, as in college admissions with affirmative action
\cite{AyTu20a,AyBo21a}, students have preferences over the categories.

Closest to our model in terms of students' preferences is controlled school choice with
soft bounds and overlapping types \cite{KHIY17a}.
There, each student can belong to several types, has preferences over types as well as
schools, and each type is subject to soft lower and upper bounds.
As in our model, a student can care about the category under which she is admitted.
However, the priorities of schools are ordinal, and no constraint that compares the scores
of admitted and rejected students, such as score-feasibility, is considered.
Rules that combine priorities with diversity policies have also been examined in terms of
explainability \cite{arnosti2024explainable}.

\subsection{Affirmative Action in India}

India implements one of the most comprehensive reserve systems
\cite{Gala84a,SoYe22a}, and its admission and recruitment procedures have been analyzed from a
market design perspective \cite{AyTu17a,BCC+19b,SoYe22a}.
The transfer of unfilled reserved seats to other categories (de-reservation) has been
formulated through choice rules with capacity transfers \cite{AyTu20a}, and a choice rule
that jointly implements reservation and de-reservation of OBC seats has been proposed,
together with a characterization of the deferred acceptance mechanism based on it
\cite{AyTu23a}.
These studies mainly design the choice rules of institutions and combine them with the
deferred acceptance or cumulative offer mechanism.
In contrast, we treat the bounded-gap requirement as a constraint on the matching
itself, and we focus on how the protection level of reserved categories interacts with
the incentives of students.

\subsection{Matching with Constraints}

Our model is also related to matching with distributional constraints
\cite{KaKo15a,GKK+17a,KTY18a}; see \cite{ABY22a} for a survey.
In these models, feasibility is usually defined by the number of students assigned to
each school. Score-feasibility, in contrast, depends on the scores and preferences of
rejected students, so it does not fit directly into these frameworks.
Our mechanisms are variants of serial dictatorship \cite{Sven99a,AbSo98a,GKK+17a}, where the order of
students is given by scores.

A recent work \cite{SZZ26} studies the interaction between overlapping eligibility, one-to-one category counting, and score-feasibility.  
It compares a student-proposing score-elevated deferred acceptance
mechanism (SPSE) with a preference-aware overlapping de-reservation
mechanism (PA-ODR), focusing on the
feasibility of filling overlapping reserve quotas.
The present paper addresses a different question:
how much reserve capacity should be protected when
de-reservation is permitted.
We introduce a rejection-based open cutoff, rather than the
lowest score of admitted open-category students, and allow
category-specific score margins.
We then study RCP and establish
trade-offs among SP, RCP, and WNW through several serial dictatorship mechanisms.

\section{\uppercase{Model}}
\label{sec:model}

A market consists of the following components.
\begin{itemize}
\item A finite set of students $I$.
\item A single college, which has the \emph{open category} $o$ and a finite set of
\emph{reserved categories} $T$. Let $V:=T\cup\{o\}$.
\item $q=(q_t)_{t\in T}$, where each $q_t\in\mathbb{Z}_{\ge 0}$ is the maximum quota of 
category $t$. 
\item The college's (total) maximum quota $Q \in\mathbb{Z}_{> 0}$. We assume $Q > \sum_{t \in T} q_t$. 
\item For each student $i\in I$, let 
$V(i)\subseteq V$ denote the categories for which $i$ is eligible. 
We assume $o\in V(i)$ for each $i\in I$.
\item The set of contracts
$X:=\{(i,c)\mid i\in I,\ c\in V(i)\cup\{\emptyset\}\}$.
Contract $(i,c)$ means that student $i$ is admitted to the college under
category $c$, and $(i,\emptyset)$ means that $i$ is assigned to the outside option,
i.e., she is not admitted under any category. 
For a set of contracts $Y\subseteq X$,
let $Y_i:=\{(j,c')\in Y \mid j=i\}$ denote the contracts of student $i$ in $Y$, 
and $Y_c:=\{(j,c')\in Y \mid c'=c\}$ denote the contracts of category $c$ in $Y$.

\item $\succ_I=(\succ_i)_{i\in I}$. 
Each $\succ_i$ is a strict preference of student $i$
over $X_i$. 
For brevity, we write $c \succ_i c'$ 
instead of $(i,c)\succ_i (i,c')$. We write $c \succeq_i c'$ 
when $c \succ_i c'$ or $c=c'$ holds. 
% We also use a notation like $Y_i \succ_i c'$, which means
% $(i,c) \succ_i (i,c')$ where $Y_i=\{(i,c)\}$.
Category $c$ is \emph{acceptable} for $i$ if $c\succ_i\emptyset$.
\item Each student $i$ has a score $\sigma(i)\ge 0$. We assume that no two students
have the same score. 
This assumption is without loss of 
generality; when two students 
have the same score, we can apply some tie-breaking rule to decide their order. 
\item For each reserved category $t\in T$, a margin $\delta_t\ge 0$ is given.
\end{itemize}
We say $Y\subseteq X$ is a \emph{matching} if, for each student $i\in I$, $|Y_i|=1$
holds and the unique contract $(i,c)\in Y_i$ satisfies $c\succeq_i\emptyset$.
We say $Y\subseteq X$ is a \emph{partial matching} if, for each student $i\in I$,
$|Y_i|\le 1$ holds and $(i,c)\in Y_i$ implies $c\succeq_i\emptyset$.
With a slight abuse of notation, when $Y_i=\{(i,c)\}$, we write $c'\succ_i Y_i$ (or $Y_i \succ_i c'$)
for $c'\succ_i c$ (or for $c \succ_i c'$). 

In the rest of this section, we introduce several necessary or desirable properties. 
One basic property is feasibility. We introduce the following two alternative 
definitions. 

\begin{definition}[Feasibility (fixed quota)]
\label{def:feasibility}
Let $q_o= Q - \sum_{t \in T} q_t$. 
A matching $Y$ is \emph{feasible} (in the fixed quota model) 
if $|Y_c|\le q_c$ holds for each $c\in V$.
\end{definition}
Here, we assume the maximum quota of the open category is fixed to $q_o$, which 
is equal to $Q - \sum_{t\in T} q_t$. 

Alternatively, we can assume the open category does not have a fixed quota; more students 
can be assigned to the open category as long as the total maximum quota $Q$ is satisfied. 
\begin{definition}[Feasibility (flexible quota)]
\label{def:flexfeasibility}
A matching $Y$ is \emph{feasible} (in the flexible quota model) 
if $|Y_t|\le q_t$ holds for each $t\in T$, and 
$\sum_{c \in V} |Y_c| \leq Q$ holds. 
\end{definition}

For a matching $Y$, let
\begin{equation}
\label{eq:Ro}
R_o(Y):=\{i\in I\mid (i,\emptyset)\in Y,\ o\succ_i\emptyset\},
\end{equation}
i.e., $R_o(Y)$ is the set of students who are assigned to the outside option although
the open category is acceptable for them.
We define $e_o(Y):=\max_{i\in R_o(Y)}\sigma(i)$ if $R_o(Y)$ is not empty, and
$e_o(Y):=0$ otherwise.
Intuitively, $e_o(Y)$ is the highest score among students who are rejected from the college
even though they want to be admitted to the open category.

\begin{definition}[Score-feasibility]
\label{def:sf}
A matching $Y$ is \emph{score-feasible} if $\sigma(i)\ge e_o(Y)-\delta_t$ holds for each
$(i,t)\in Y$ with $t\in T$.
\end{definition}

Score-feasibility bounds the gap between the cutoff score of the open category and the
scores of the students admitted through reserved categories.
The margin $\delta_t$ is defined for each category; different categories can have different margins.
When $\delta_t$ is sufficiently large for every $t\in T$, the requirement becomes vacuous.
When $\delta_t=0$, a student can be admitted under $t$ only if her score is higher than or equal to 
the score of every student in $R_o(Y)$.

Feasibility and score-feasibility can be applied to a partial matching $Y$, 
by assuming $i$ is assigned to the outside option $\emptyset$ if $Y_i=\emptyset$. 

An alternative way to define score-feasibility is to use the lowest score of the students
admitted under $o$, rather than the highest score of the rejected students.
The two definitions are similar, and their difference would be small in practice when there are
many borderline students.
We apply the current definition since the following property holds. 
This property is crucial to guarantee that our proposed mechanisms satisfy 
several desirable properties.
\begin{proposition}
\label{prop:remove}
Assume that a matching $Y$ is score-feasible, and let $(j,c)\in Y$ with
$c\neq\emptyset$ be the contract of the admitted student with the lowest score.
Then the (partial) matching $(Y\setminus\{(j,c)\})$ is also
score-feasible.
\end{proposition}
\begin{proof}
Let $Y':=(Y\setminus\{(j,c)\})$.
When $o \succ_j \emptyset$, $R_o(Y') = R_o(Y) \cup \{j\}$. Otherwise, $R_o(Y') = R_o(Y)$. 
In both cases, either (i) $e_o(Y')=e_o(Y)$ or (ii) $e_o(Y')=\sigma(j)$ holds. 
In case (i), since $Y$ is score-feasible, for each $(i, t) \in Y'$, where $t\in T$, 
$\sigma(i) \ge e_o(Y) - \delta_t = e_o(Y') - \delta_t$ holds. Thus, $Y'$ is score-feasible. 
In case (ii), since $\sigma(j) < \sigma(i)$ for each $(i, t) \in Y'$, where
$t \in T$, 
$\sigma(i) > \sigma(j) \geq e_o(Y') - \delta_t$ 
% \zh{Here when c is the open category, $\delta_c$ is not defined.}
holds. 
Thus, $Y'$ is score-feasible. 
\end{proof}
If $Y$ is feasible and $Y' \subset Y$, $Y'$ is also feasible in both 
fixed quota and flexible quota models. 
Thus, removing the lowest-score admitted student from a feasible and score-feasible
matching keeps both properties.
This property does not hold under the alternative definition.
Let $T=\{t\}$ and $\delta_t=10$, $q_t=1$, $Q=3$, and consider three students with
$\sigma(i_1)=100$, $V(i_1)=\{o\}$; $\sigma(i_2)=80$, $V(i_2)=\{o,t\}$; and
$\sigma(i_3)=60$, $V(i_3)=\{o\}$.
The matching $Y=\{(i_1,o),(i_2,t),(i_3,o)\}$ is feasible and score-feasible under both
definitions.
Removing the lowest-score student $i_3$ yields $Y'=\{(i_1,o),(i_2,t)\}$, where
we assume $i_3$ is assigned to the outside option. 
$Y'$ is score-feasible under our definition, since every admitted student has a higher
score than $i_3$.
Under the alternative definition, however, $Y'$ is not score-feasible: the lowest score in
$o$ is now $100$, while $\sigma(i_2)=80<100-\delta_t$.

Fairness requires that if two students compete for a seat in a category, 
it is resolved based on their scores; the student with a higher score obtains the seat.
\begin{definition}[Fairness]
\label{def:fair}
Student $i$ has \emph{justified envy} toward student $j$ in matching $Y$ if
$Y_j=\{(j,c)\}$, $c\in V(i)$, $c\succ_i Y_i$, and $\sigma(i)>\sigma(j)$ hold.
A matching is \emph{fair} if no student has justified envy.
\end{definition}

Non-wastefulness (NW) requires that we cannot move a student to a better category while keeping 
the other assignments without violating feasibility or score-feasibility. 
\begin{definition}[Non-wastefulness (NW)]
\label{def:nw}
In matching $Y$, student $i$ \emph{claims an empty seat} of category $c$ if $c\in V(i)$,
$c\succ_i Y_i$, and the matching $(Y\setminus Y_i)\cup\{(i,c)\}$ is feasible and
score-feasible.
A matching is \emph{non-wasteful} (NW) if no student claims an empty seat.
\end{definition}

For two matchings $Y$ and $Z$, we say $Z$ \emph{Pareto dominates} $Y$ if every student
weakly prefers her contract in $Z$ to her contract in $Y$, and at least one student
strictly prefers her contract in $Z$.

\begin{definition}[Pareto efficiency (PE)]
\label{def:pe}
A feasible and score-feasible matching is \emph{Pareto efficient} (PE) if it is not
Pareto dominated by another feasible and score-feasible matching.
\end{definition}

PE implies NW: if student $i$ claims an empty seat of $c$ in $Y$, the
matching $(Y\setminus Y_i)\cup\{(i,c)\}$ is feasible, score-feasible, and Pareto
dominates $Y$.

A \emph{mechanism} is a function that maps each profile of preferences reported by the
students to a matching.
We say a mechanism satisfies a property of matchings (e.g., fairness) if it always
returns a matching that satisfies the property with respect to the reported preferences.

\begin{definition}[Strategy-proofness (SP)]
\label{def:sp}
A mechanism is \emph{strategy-proof} (SP) if no student can obtain a strictly better
contract (with respect to her true preference) by misreporting her preference, regardless
of the reports of the other students.
\end{definition}

\section{\uppercase{Fixed Quota}}
\label{sec:fixed}
We first consider the basic model, where the quota of the open category is 
fixed (Definition~\ref{def:feasibility}).
NW and PE are defined using this feasibility definition.

In this model, we can use the following mechanism. 
\begin{definition}[Score-based serial dictatorship mechanism (SSD)]
\label{def:ssd}
The \emph{score-based serial dictatorship} mechanism (SSD) works as follows.
Initially, the partial matching $Y$ contains no contract and the threshold $e$ is $0$.
Students are chosen one by one in decreasing order of score.
When student $i$ is chosen, category $c\in V(i)$ is \emph{available} for $i$ if:
\begin{enumerate}
\item $|Y_c|<q_c$ holds, and
\item if $c\in T$, $\sigma(i)\ge e-\delta_c$ holds.
\end{enumerate}
The outside option is always available.
Student $i$ is assigned to her most preferred available option according to her reported
preference, and the corresponding contract is added to $Y$.
If $i$ is the first student who is assigned to the outside option although $o$ is
acceptable for her, we set $e:=\sigma(i)$.
\end{definition}
This mechanism is an instance of the serial dictatorship mechanism~\cite{Sven99a,AbSo98a,GKK+17a}, 
where each student chooses her most preferred option among those that keep
the partial matching feasible and score-feasible.

\begin{theorem}
\label{thm:ssd}
SSD satisfies feasibility, score-feasibility, fairness, PE, and SP.
\end{theorem}
\begin{proof}
Let $Y$ be the matching obtained by SSD.
For student $i$, let $Y^i$ denote the partial matching just before $i$ is chosen.

We first examine the threshold $e$.
Since students are chosen in decreasing order of score, the first student who is
assigned to the outside option, although $o$ is acceptable for her, has the highest score
in $R_o(Y)$.
Thus, once $e$ is set, it equals $e_o(Y)$ and does not change afterward.

\textbf{Feasibility.} Clear from the construction.

\textbf{Score-feasibility.}
Assume $(i,t)\in Y$ for $t\in T$.
If $e$ is already set when $i$ is chosen, $\sigma(i)\ge e-\delta_t=e_o(Y)-\delta_t$ holds
by construction.
Otherwise, all students in $R_o(Y)$ are chosen after $i$, so
$e_o(Y)\leq \sigma(i)\le\sigma(i)+\delta_t$ holds.

\textbf{Fairness.}
Assume $i$ has justified envy toward $j$, where $Y_j=\{(j,c)\}$.
Since $\sigma(i)>\sigma(j)$, $i$ is chosen before $j$.
Since $c\succ_i Y_i$, category $c$ was not available for $i$.
There are two cases.
\begin{enumerate}
\item $c$ is full when $i$ is chosen.
Then $c$ is also full when $j$ is chosen, so $j$ cannot be assigned to $c$.
\item $c\in T$ and $\sigma(i)<e-\delta_c$ holds when $i$ is chosen.
Then $e$ is already set, so the threshold is the same when $j$ is chosen.
Since $\sigma(j)<\sigma(i)<e-\delta_c$, $j$ cannot be assigned to $c$.
\end{enumerate}
Both cases contradict $Y_j=\{(j,c)\}$.

\textbf{PE.}
Assume matching $Z$ is feasible and score-feasible, and $Z$ Pareto dominates $Y$.
Let $i$ be the student with the highest score among the students whose contracts are
different in $Y$ and $Z$.
Then $i$ strictly prefers her contract in $Z$. Let $Z_i=\{(i,c)\}$.
Since $Y^i \cup \{(i,c)\} \subseteq Z$, 
$Y^i \cup \{(i,c)\}$ is feasible. Also, since it is obtained from $Z$ by repeatedly 
removing the lowest-score student, by 
Proposition~\ref{prop:remove}, $Y^i \cup \{(i,c)\}$ is score-feasible. 
Then, $c$ must be available for $i$ in SSD. 
This contradicts the assumption that $i$ is not assigned to $c$ in SSD. 

\textbf{SP.}
The set of options available for $i$ is determined by $Y^i$ and the value of $e$ when $i$
is chosen.
They depend only on the reports of the students whose scores are higher than $\sigma(i)$.
Thus, $i$ cannot change the set of available options by her report, and reporting
truthfully is optimal for $i$.
\end{proof}

\section{\uppercase{Flexible Quota: Model and Impossibility}}
\label{sec:flexible}
In the fixed quota model, there is a chance that 
the number of students who want to be assigned to a particular reserved 
category $t$ is less than $q_t$. As a result, students who hope to be assigned 
to the open category are rejected, even though the college has vacant seats. 

The alternative feasibility definition (Definition~\ref{def:flexfeasibility}) 
is intended to avoid such a situation, where the open category 
does not have a fixed quota; the number of students assigned to 
the open category is determined by the total number of students assigned 
to reserved categories. 
In the rest of this paper, we discuss this flexible quota model. 
NW and PE are defined using this feasibility requirement. 

By using Definition~\ref{def:flexfeasibility}, the college can fully utilize
its available seats, since the college can transfer seats from a reserved category to the open category. However, there is a chance that most of the reserved category seats are transferred to the open category and very few students are admitted under 
reserved categories.

In summary, de-reservation improves seat utilization but may undermine the
purpose of reserve quotas: a reserved seat may be transferred to
the open category even when students explicitly rank that reserved
category as their first choice.
Conversely, protecting all reserved seats regardless of demand
may result in unnecessary vacancies.

To balance these concerns, we introduce
\emph{reserved category protection} (RCP) as follows. 
\begin{definition}[Reserved category protection (RCP)]
\label{def:rcp}
For each student $i$, let $\topc_i$ denote her most preferred element of
$V(i)\cup\{\emptyset\}$.
For each category $t\in T$, let
$d_t:=\min\{q_t,\ |\{i\in I\mid\topc_i=t\}|\}$.
A matching $Y$ satisfies \emph{reserved category protection} (RCP) if
\begin{equation}
\label{eq:rcpdef}
|Y_o|\le Q- \sum_{t\in T}\bigl(\max\{|Y_t|,d_t\}\bigr)
\end{equation}
holds.
\end{definition}
Intuitively, we protect $d_t$ seats for each reserved category $t$, corresponding to the reported first-choice demand.
A reserved seat may be transferred to the open category only if
it is neither occupied nor protected by such demand.
Note that RCP protects capacity rather than guaranteeing admission:
some protected seats may remain unfilled because of
score-feasibility constraints.

\begin{comment}
To avoid such a situation, we set a safeguard that restricts the number of 
transferred seats from each reserved category to the open category, 
which we call \emph{reserved category protection}. 

% \zh{It only checks the number of students who rank $t$ as top category, but not whether it violates score-feasibility when matched.}

Here, $d_t$ is the number of seats of category $t$ that are demanded as the top choice.
RCP means that a seat of $t$ can be transferred to the open category
only if it is neither used nor demanded in this sense.
% Note that (\ref{eq:rcpdef}) implies $\sum_{c\in V}|Y_c|\le Q$.
\end{comment}

RCP is a rather weak requirement; it only considers the top category of each student. 
However, it turns out that this weak requirement becomes incompatible 
with NW, as shown in the following theorem. 
\begin{theorem}
\label{thm:imp1}
There is a case where no matching simultaneously satisfies feasibility,
score-feasibility, NW, and RCP.
\end{theorem}
\begin{proof}
Consider the case where $T=\{t\}$, $q_t=1$, $Q=2$, and $\delta_t=10$.
There are three students:
\begin{itemize}
\item $i_1$: $\sigma(i_1)=100$, $V(i_1)=\{o\}$, $o\succ_{i_1}\emptyset$,
\item $i_2$: $\sigma(i_2)=90$, $V(i_2)=\{o\}$, $o\succ_{i_2}\emptyset$,
\item $i_3$: $\sigma(i_3)=60$, $V(i_3)=\{o,t\}$,
$t\succ_{i_3}o\succ_{i_3}\emptyset$.
\end{itemize}
Assume matching $Y$ satisfies feasibility, score-feasibility, and RCP.
Since $d_t=1$, RCP requires $|Y_o|\le 1$.
Thus, at least one of $i_1$ and $i_2$ is assigned to the outside option, and
$e_o(Y)\ge 90$ holds.
Then, assigning $i_3$ to $t$ violates score-feasibility, since $60<90-\delta_t$.
Therefore, no student is assigned to $t$ and at most one student is assigned to $o$.
Then, at least one of $i_1$ and $i_2$ is assigned to the outside option, and
moving her to $o$ satisfies feasibility (at most two students are assigned in total) and
score-feasibility (no student is assigned to $t$).
Thus, she claims an empty seat of $o$; $Y$ violates NW.
\end{proof}

Since NW is too strong, we introduce a weaker efficiency requirement as follows. 
\begin{definition}[Weak non-wastefulness (WNW)]
\label{def:wnw}
In matching $Y$, student $i$ \emph{strongly claims an empty seat} of category $c$ if
$c\in V(i)$, $c\succ_i Y_i$, and the matching $(Y\setminus Y_i)\cup\{(i,c)\}$ satisfies
feasibility, score-feasibility, and RCP.
A matching is \emph{weakly non-wasteful} (WNW) if no student strongly claims an empty
seat.
\end{definition}

NW implies WNW, since a student who strongly claims an empty seat also
claims it.

We can also weaken Pareto efficiency based on RCP. 
\begin{definition}[Constrained Pareto efficiency]
\label{def:cpe}
A matching that satisfies feasibility, score-feasibility, and RCP is
\emph{constrained Pareto efficient} if it is not Pareto dominated by another matching
that satisfies these three conditions.
\end{definition}
As in the case of PE and NW, constrained Pareto efficiency implies WNW:
if $i$ strongly claims an empty seat of $c$ in $Y$, the matching
$(Y\setminus Y_i)\cup\{(i,c)\}$ satisfies the three conditions and Pareto dominates $Y$.
Note that even when a matching $Y$ satisfies PE,
$Y$ is not necessarily constrained Pareto efficient, 
since PE does not require RCP. 

Even under these weaker requirements, an impossibility result remains if we also
require fairness and SP.

\begin{theorem}
\label{thm:imp2}
No mechanism simultaneously satisfies feasibility, score-feasibility, fairness, WNW,
RCP, and SP in general.
\end{theorem}
\begin{proof}
Consider the case where $T=\{t\}$, $q_t=1$, $Q=2$, and $\delta_t=10$.
There are three students:
\begin{itemize}
\item $i_1$: $\sigma(i_1)=100$, $V(i_1)=\{o\}$, $o\succ_{i_1}\emptyset$,
\item $i_2$: $\sigma(i_2)=90$, $V(i_2)=\{o\}$, $o\succ_{i_2}\emptyset$,
\item $i_3$: $\sigma(i_3)=80$, $V(i_3)=\{o,t\}$,
$o\succ_{i_3}t\succ_{i_3}\emptyset$.
\end{itemize}
Assume a mechanism satisfies all the properties except SP, and let $Y$
be the matching it returns.

\textbf{Truthful reports.}
Since $d_t=0$, no seat for $t$ is protected. 
At most two students can be assigned in total. 
If $i_3$ is assigned to $o$, 
at least one of $i_1$ and $i_2$ is assigned to the outside option,
and she has justified envy toward $i_3$.
Similarly, if $i_3$ is assigned to $t$, 
score-feasibility requires $e_o(Y)\le 90$, so $i_1$ is
assigned to $o$, and $i_2$ is assigned to the outside option.
Then, moving $i_3$ to $o$ satisfies feasibility, score-feasibility (no student is
assigned to $t$), and RCP.
Thus, $i_3$ strongly claims an empty seat of $o$.
Therefore, $i_3$ is assigned to the outside option.
The only matching satisfying the required properties is to 
assign both $i_1$ and $i_2$ to $o$. 

\textbf{Misreport.}
Assume $i_3$ reports $t\succ'_{i_3}o\succ'_{i_3}\emptyset$.
Since $d_t=1$, RCP requires $|Y_o|\le 1$.
If $i_2$ or $i_3$ is assigned to $o$, $i_1$ is assigned to the outside option and has
justified envy.
If no student is assigned to $o$, moving $i_1$ to $o$ satisfies feasibility and RCP.
It also satisfies score-feasibility: even if $i_3$ is assigned to
$t$, the rejected student with the highest score is $i_2$ and $80\ge 90-\delta_t$ holds.
Thus, $i_1$ strongly claims an empty seat of $o$.
Therefore, $i_1$ is assigned to $o$, $i_2$ is assigned to the outside option, and
$e_o(Y)=90$.
If $i_3$ is assigned to the outside option, moving $i_3$ to $t$ satisfies feasibility,
score-feasibility ($80\ge 90-\delta_t$), and RCP, so $i_3$
strongly claims an empty seat of $t$ with respect to the reported preference.
Therefore, $i_3$ is assigned to $t$.

Since $t\succ_{i_3}\emptyset$ holds in her true preference, $i_3$ obtains a strictly
better contract by misreporting. Thus, SP is violated.
\end{proof}


\section{\uppercase{Mechanisms and Tight Boundary}}
\label{sec:boundary}
By Theorem~\ref{thm:imp2}, assuming feasibility, score-feasibility, and fairness are
indispensable, we must give up at least one of SP, RCP, and WNW.
In this section, we show that sacrificing any one of them is enough;
we present a mechanism that satisfies any pair of these properties. 

\subsection{RCP-SSD: Sacrificing SP}
\begin{comment}
\zh{When assigning a student to a reserved category $t$,
the RCP-SSD algorithm checks
\[
|Y_t|<q_t,\qquad
\sum_{c\in V}|Y_c|<Q,\qquad
\sigma(i)\ge e-\delta_t.
\]
However, these conditions do not guarantee that the
resulting assignment preserves RCP.
\\
Consider an instance with
\[
T=\{a,b\},\quad Q=3,\quad
q_a=q_b=1,\quad
\delta_a=\delta_b=10.
\]
There are four students with distinct scores:
\[
\sigma(i_1)=100,\quad
\sigma(i_2)=90,\quad
\sigma(i_3)=80,\quad
\sigma(i_4)=70.
\]
Their preferences are
\[
\begin{aligned}
&i_1,i_2:\quad o\succ\varnothing,\\
&i_3:\quad o\succ b\succ\varnothing,\\
&i_4:\quad a\succ o\succ\varnothing.
\end{aligned}
\]
Thus, $d_a=1$ and $d_b=0$.
The algorithm first assigns $i_1$ and $i_2$ to $o$.
When considering $i_3$, the open category is unavailable
because the remaining seat is protected for $a$.
However, category $b$ is available under the current
definition, since its quota is unfilled, the total number
of admitted students is less than $Q$, and the score
condition is satisfied.
Hence, $i_3$ is assigned to $b$.
The resulting partial matching violates RCP:
\[
|Y_o|+\max\{|Y_a|,d_a\}
+\max\{|Y_b|,d_b\}
=2+1+1=4>Q.
\]
Therefore, the current availability conditions do not
maintain RCP, contrary to the argument in Theorem~4.}
\end{comment}

\begin{definition}[Reserved category protecting SSD (RCP-SSD)]
\label{def:rcpssd}
The \emph{reserved category protecting SSD} mechanism (RCP-SSD) first computes $d_t$ for each $t\in T$ from the reported preferences.
Then it runs SSD, where the availability is modified as follows.
When student $i$ is chosen at partial matching $Y$, 
\begin{itemize}
    \item $t\in V(i)\cap T$ is available 
if $|Y_t|< q_t$, $|Y_o| + \max\{|Y_t|+1, d_t\} + \sum_{s \in T\setminus\{t\}}\max\{|Y_s|, d_s\}\leq Q$, 
and $\sigma(i)\ge e-\delta_t$. 
\item $o$ is available 
if $|Y_o|+1 + \sum_{t \in T}\max\{|Y_t|, d_t\}\leq Q$.
\end{itemize}
The threshold $e$ is maintained in the same way as SSD.
\end{definition}
In words, student $i$ can be assigned to category $c$ as long as feasibility, score-feasibility, and RCP are satisfied. 

\begin{comment}
\zh{I suggest the following modification:
\\
For a reserved category $t\in T$, the category is
available to student $i$ if
\[
|Y_t|<q_t,
\qquad
\sigma(i)\ge e-\delta_t,
\]
and
\[
|Y_o|
+\max\{|Y_t|+1,d_t\}
+\sum_{s\in T\setminus\{t\}}
\max\{|Y_s|,d_s\}
\le Q.
\]
The open category is available if
\[
|Y_o|+1+
\sum_{t\in T}\max\{|Y_t|,d_t\}
\le Q.
\]
In particular, a student assigned to a reserved category
may use an already protected seat, but cannot consume
capacity protected for another category.
}
\end{comment}
\begin{theorem}
\label{thm:rcpssd}
RCP-SSD satisfies feasibility, score-feasibility, fairness, RCP, and 
constrained Pareto efficiency (thus, it also satisfies WNW). 
\end{theorem}
\begin{proof}
\textbf{Feasibility and RCP.}
The mechanism keeps the condition (\ref{eq:rcpdef})
during the execution.
Thus, the final matching satisfies RCP. 
Also, $d_t\leq q_t$ by definition. 
$|Y_t|\leq q_t$ holds for each $t \in T$ during the execution.
The mechanism keeps RCP, i.e., 
$|Y_o| + \sum_{t\in T} \bigl(\max\{|Y_t|,d_t\}\bigr)\le Q$. 
Since $|Y_t| \leq \max\{|Y_t|,d_t\}$, 
$|Y_o| + \sum_{t\in T} |Y_t| \le Q$ holds. 
Thus, the final matching satisfies feasibility. 

\textbf{Score-feasibility.}
The proof is the same as that of Theorem~\ref{thm:ssd}.

\textbf{Fairness.}
Assume $c$ is not available for $i$. Since $e$ weakly monotonically increases, and the number of assigned students
to each category weakly monotonically increases, 
$c$ is not available for the students assigned after $i$, whose scores are lower than $\sigma(i)$.

\textbf{Constrained Pareto efficiency.}
Assume a matching $Z$ satisfies feasibility, score-feasibility, and RCP, and $Z$ Pareto
dominates $Y$. Let $i$ be the student with the highest score 
among the students whose contracts are different in $Y$ and 
$Z$. Assume $i$ is assigned to $c$ in $Z$. 
Let $Y^i$ denote the partial matching 
obtained just before assigning $i$ in the mechanism. 
Since $Y^i\cup\{(i,c)\}\subseteq Z$, 
$Y^i\cup\{(i,c)\}$ is feasible and satisfies RCP. 
Also, 
since it is obtained from $Z$ by repeatedly removing 
the lowest-score student, by Proposition~\ref{prop:remove},
it is score-feasible. Then, $c$ must be available for 
$i$ in RCP-SSD. This contradicts the assumption that
$i$ is not assigned to $c$ in RCP-SSD.

\textbf{WNW.}
Constrained Pareto efficiency implies WNW.
\end{proof}

By Theorem~\ref{thm:imp2}, RCP-SSD violates SP.
Indeed, in the case used in the proof of Theorem~\ref{thm:imp2}, RCP-SSD assigns
$i_1$ and $i_2$ to $o$ and $i_3$ to the outside option under truthful reports.
When $i_3$ reports $t$ as her top choice, $d_t$ changes from $0$ to $1$, only one student
can be assigned to $o$, and $i_3$ is assigned to $t$.
That is, a student can manipulate the protection level through her report.
Also, by Theorem~\ref{thm:imp1}, RCP-SSD violates 
NW, and hence PE: in the case used in the proof of Theorem~\ref{thm:imp1}, it assigns only $i_1$
to the college, and the seat protected for $t$ remains empty.




\subsection{Flexible SSD: Sacrificing RCP}

\begin{definition}[Flexible SSD]
\label{def:flexssd}
The \emph{flexible SSD} mechanism runs SSD, where the availability is
modified as follows.
When student $i$ is chosen at partial matching $Y$, 
\begin{itemize}
    \item $t\in V(i)\cap T$ is available if $|Y_t|< q_t$ and 
    $\sum_{c \in V} |Y_c| < Q$.
\item $o$ is available if $\sum_{c \in V} |Y_c| < Q$.
\end{itemize}
\end{definition}
Note that although 
flexible SSD does not maintain the threshold $e$, 
it satisfies score-feasibility as shown in Theorem~\ref{thm:flexible}.



\begin{theorem}
\label{thm:flexible}
Flexible SSD satisfies feasibility, score-feasibility, fairness, PE, and
SP. Hence, it also satisfies NW and WNW.
\end{theorem}
\begin{proof}
Feasibility is clear from the construction. 
Fairness is also clear since if student $i$ cannot be 
assigned to $c \in V(i)$, then $c$ is not available for
the students assigned after $i$. 
SP and PE are derived by using the same argument as in the proof of Theorem~\ref{thm:ssd}.
PE implies NW, which in turn implies
WNW.

For score-feasibility, 
the open category $o$ is available as long as the total number 
of students assigned to the college is less than $Q$. 
Thus, if student $i$ is rejected from the college although she hopes to be assigned to $o$, the college has already accepted $Q$ students. 
Thus, no more students are assigned. 
As a result, score-feasibility is satisfied. 
\end{proof}

Flexible SSD does not satisfy RCP in general.
In the case used in the proof of Theorem~\ref{thm:imp1}, flexible SSD assigns $i_1$ and
$i_2$ to $o$, although $d_t=1$ requires $|Y_o|\le 1$.

\subsection{Fixed Quota SSD: Sacrificing WNW}

The third mechanism is the simplest one: it ignores the flexibility of quotas and
uses SSD of the fixed quota model as it is, 
where $q_o$ is set to $Q - \sum_{t\in T} q_t$.

\begin{definition}[Fixed quota SSD]
\label{def:fqssd}
The \emph{fixed quota SSD} mechanism runs SSD of the fixed quota model
(Definition~\ref{def:ssd}), where the quota of $o$ is $q_o=Q-\sum_{t\in T}q_t$ and the
quota of each $t\in T$ is $q_t$.
\end{definition}

The key observation is that the protection level used in fixed quota SSD is $q_t$,
which does not depend on the reports of the students, and $q_t$ is at least $d_t$.

\begin{lemma}
\label{lem:fixedrcp}
Assume for a matching $Y$, 
$|Y_t|\le q_t$ holds for each $t\in T$,
and $|Y_o| \leq Q - \sum_{t \in T} q_t$ holds. 
Then, $Y$ is feasible (in the flexible quota model) 
and satisfies RCP.
\end{lemma}
\begin{proof}
The feasibility requires that (i) $|Y_t| \leq q_t$ for each 
$t \in T$, and (ii) $\sum_{c \in V} |Y_c| \leq Q$. 
(i) is satisfied by our assumption. 
(ii) is also satisfied since 
$\sum_{c \in V} |Y_c| = \sum_{t \in T} |Y_t| + |Y_o| 
\leq Q$. 

For RCP, for each $t\in T$, $|Y_t|\leq q_t$ holds. 
Since $d_t \leq q_t$, $\max\{|Y_t|,d_t\} \leq q_t$ holds
for each $t \in T$. 
Thus, $|Y_o| \le Q - \sum_{t \in T} q_t 
\le  Q - \sum_{t \in T} \max\{|Y_t|,d_t\}$ holds. 
\end{proof}

\begin{theorem}
\label{thm:fqssd}
Fixed quota SSD satisfies feasibility, score-feasibility, fairness, RCP, and SP.
\end{theorem}
\begin{proof}
Fairness and score-feasibility are derived using the same 
argument as in the proof of Theorem~\ref{thm:ssd}.
Also, since $q_o$ is determined independently of 
students' preferences, SP is also satisfied by the same 
argument as in the proof of Theorem~\ref{thm:ssd}.
By the construction, for the obtained matching $Y$, 
$|Y_t|\le q_t$ holds for each $t\in T$,
and $|Y_o| \leq Q - \sum_{t \in T} q_t$ holds. 
Then, by Lemma~\ref{lem:fixedrcp},
$Y$ satisfies feasibility and RCP. 
\end{proof}

In contrast to RCP-SSD, fixed quota SSD protects every reserved seat, whether it is
demanded or not.
This protection is excessive in some cases.
Thus, WNW fails.

\begin{example}
\label{ex:fqssd}
Consider the case in the proof of Theorem~\ref{thm:imp2} with truthful reports, where
$q_t=1$, $Q=2$, and $\delta_t=10$. Then, fixed quota
SSD sets $q_o$ to $1$. 
First, $i_1$ is assigned to $o$.
Next, $i_2$ is assigned to the outside option since $o$ is full, and $e=90$ is set.
Finally, $i_3$ is assigned to $t$, since $o$ is full and $80\ge 90-\delta_t$ holds.
The outcome is the same when $i_3$ reports $t\succ'_{i_3}o\succ'_{i_3}\emptyset$, so
$i_3$ cannot obtain a better contract by misreporting.

However, $i_3$ prefers $o$ to $t$.
Let $Z$ be the matching obtained by moving $i_3$ from $t$ to $o$.
Then $|Z_o|=2$ and $|Z_t|=0$, so $Z$ is feasible.
It is score-feasible since no student is assigned to $t$.
Since $d_t=0$, one seat can be transferred from $t$ to $o$;
$Z$ satisfies RCP.
Thus, $i_3$ strongly claims an empty seat of $o$, and the outcome of fixed quota SSD violates WNW.
\end{example}

By Theorem~\ref{thm:imp2}, the failure of WNW is unavoidable for any mechanism that
satisfies feasibility, score-feasibility, fairness, RCP, and SP.

\subsection{Tight Boundary}

\begin{table*}[t]
\caption{Properties of the three mechanisms in the flexible quota model.}\label{tab:comparison} \centering
\begin{tabular}{|l|c|c|c|}
  \hline
   & RCP-SSD & Flexible SSD & Fixed quota SSD \\
  \hline
  Feasibility & \checkmark & \checkmark & \checkmark \\
  \hline
  Score-feasibility & \checkmark & \checkmark & \checkmark \\
  \hline
  Fairness & \checkmark & \checkmark & \checkmark \\
  \hline
  \hline
  Strategy-proofness (SP) & $\times$ & \checkmark & \checkmark \\
  \hline
  Reserved category protection (RCP) & \checkmark & $\times$ & \checkmark \\
  \hline
  Weak non-wastefulness (WNW) & \checkmark & \checkmark & $\times$ \\
  \hline
  \hline
  Non-wastefulness (NW) & $\times$ & \checkmark & $\times$ \\
  \hline
  Pareto efficiency (PE) & $\times$ & \checkmark & $\times$ \\
  \hline
  Constrained Pareto efficiency & \checkmark & $\times$ & $\times$ \\
  \hline
\end{tabular}
\end{table*}

Table~\ref{tab:comparison} summarizes the properties of the three mechanisms.
All entries follow from Theorems~\ref{thm:imp1}--\ref{thm:fqssd}, Example~\ref{ex:fqssd}, and the remarks following these results.
Thus, we obtain 
a tight boundary between possibility and
impossibility: the impossibility of Theorem~\ref{thm:imp2} disappears as soon as any one
of SP, RCP, or WNW is dropped.

The three mechanisms can also be understood in a unified way in terms of the protection
level of each reserved category $t$, i.e., the number of seats of $t$ that cannot be transferred to
the open category.
Flexible SSD has no protection, which can be considered as level~$0$, 
RCP-SSD can be considered level~$d_t$ where $d_t$ seats are protected, 
and fixed quota SSD can be considered level~$q_t$, where all seats are protected. 
A protection level of $0$ violates RCP, and a protection level of $q_t$ is too
conservative to satisfy WNW.
The intermediate level $d_t$ satisfies both, but it depends on students' reports,
and this dependence is the source of manipulation.

\section{\uppercase{Conclusions}}
\label{sec:conclusion}

We studied reserve systems in which the scores of the students admitted through reserved
categories cannot be too far below the cutoff score of the open category, where the
allowed gap can differ across reserved categories.
For fixed quotas, SSD satisfies all the desirable
properties.
For flexible quotas, where unfilled reserved seats can be used for the open category,
protecting the reserved seats demanded by students is incompatible with
NW. 
Moreover, when feasibility, score-feasibility, and fairness are required,
SP, RCP, and WNW cannot be
satisfied simultaneously, while each pair of them is achieved by RCP-SSD, flexible SSD,
or fixed quota SSD.
Hence, our results give a tight boundary between possibility and impossibility.

There are several directions for future work.
First, we considered a single college. Extending the model to many colleges, where a
student has preferences over college--category pairs, is an important next step.
The difficulty is that score-feasibility depends on which students are rejected.
Second, we assumed that each student is assigned to at most one category.
% and that categories are not nested. 
%%% MY. We have no assumption on categories, so it is possible that categories can be nested; it does not matter since a student is counted at most one category. 
Handling overlapping reservations such as the vertical and
horizontal reservations in India \cite{SoYe22a} is left for future work.
Third, protection levels between $d_t$ and $q_t$ that are independent of students'
reports, such as levels estimated from past data, are worth investigating.
Finally, it would be interesting to quantify the efficiency loss of fixed quota SSD and
the incentive to manipulate RCP-SSD through simulations with realistic data.


\section*{\uppercase{Acknowledgments}}

\begin{itemize}
  \item \textbf{AI Tools:} 
  The authors declare the use of Generative Artificial Intelligence tools during the preparation of this manuscript. Specifically, Claude and ChatGPT were used for language proofreading and grammatical refinement. After using these tools, the authors reviewed and edited the content as needed and take full responsibility for the content of the publication. 
\end{itemize}


\bibliographystyle{apalike}
{\small
\bibliography{references}}

\end{document}
